\clearpage
\subsection{Both finite tails are canonical}
\label{sec:finite-tails}

The finite tails can be compared in finite type~A before applying the
common affine prefix. In the finite spherical module $C_V\mathcal H(S_N)$ put
\[
 \mathsf a_m=\frac{C_VC_KC_V}{\pi_L^2},\qquad
 \mathsf b_m=\frac{C_VC_{K'}C_E}
 {\pi_{V\cap K'}\pi_{K'\cap E}}.
\]
The right $V$-eigenrelations of $F=C_IC_J/\pi_V$ define the module map
\[
 \Phi:C_V\mathcal H(S_N)\longrightarrow C_I\mathcal H(W),
 \qquad \Phi(C_Vh)=Fh.
\]
Thus $\Phi(\mathsf a_m)=\mathcal R_m$ and
$\Phi(\mathsf b_m)=\mathcal{BS}_m$, with all normalization factors included.
Both finite tails are single ordinary canonical elements:
\[
 \mathsf a_m=C_{w_A},\qquad
 w_A=w_0w_{\{s_m,s_{m+1}\}},
 \qquad
 \mathsf b_m=C_{w_B},\qquad
 w_B=[m,m+1,N,N-1,\ldots,m+2,m-1,\ldots,1].
\]
Their minimal left $V$-representatives have lengths $m^2+2m-3$
and $m^2-1$, respectively.

\noindent\textit{Proof.}
In the three-letter notation of the parabolic paper, let
$c_0=c_{\mathsf A^m\mathsf B\mathsf C^m}$, with window
$[1-m,\ldots,0,m+1,N+1,\ldots,N+m]$ and length $m^2$.
For $x\in{}^VW_I$, the map
$C_VH_x\mapsto M_{c_0x}$ identifies the finite standard basis with
the type-$m$ standard span. The window formula gives
$\ell(c_0x)=m^2+\ell(x)$, and the equal-letter and unequal-letter
multiplication rules show that it commutes with finite right Hecke
actions (B, Lemmas 3.8--3.9).
The type-$m$ part of $F$ is just $M_{c_0}$ with coefficient one
(B, Lemma 3.19).
Consequently the standard coefficients of the two finite tails are
exactly the type-$m$ coefficients of the two affine products.

At the boundary $m-r=0$, our six-family estimates give coefficient
one at the unique top and strictly positive valuation at every
proper orbit maximum; refinement adds the nonnegative length
deficit (A, Lemmas 4C.2 and 4D.1).
The lowest-exponent table in the parabolic paper at type $m$ has a unique $AB$ orbit
maximum of valuation zero; the $AA$ maxima have valuation one,
and every nonmaximum has an additional positive length deficit
(B, Proposition 3.21). The top coefficient is one in each case.
The two top words are $\mathsf C^{m-1}\mathsf A\mathsf B\mathsf C
\mathsf A^{m-1}$ and $\mathsf A\mathsf B\mathsf C^m\mathsf A^{m-1}$,
which give the displayed finite permutations. Integrality and bar
invariance of the normalized products now identify both tails by
canonical triangularity.\hfill$\square$

Let $L_m=C_{s_2}\cdots C_{s_m}C_{s_1}\cdots C_{s_{m-1}}$.
The corresponding reflection word takes $w_B$ to $w_A$, increasing
length at every step: first it moves $m+1$ past the larger entries,
then it moves $m$ past them. Finite KL multiplication therefore gives
\begin{align*}
 \mathsf b_mL_m
 &=\mathsf a_m+\sum_{z<w_A}d_z(v)C_z,
 &d_z(v)&\in\mathbb N[v,v^{-1}],\\
 \mathcal{BS}_mL_m
 &=\mathcal R_m+\sum_{z<w_A}d_z(v)\Phi(C_z).
\end{align*}
Only finite canonical indices with left $V$-descents occur in the sum.
The top multiplicity is one. This expresses our web as the image of
the top finite canonical summand after the two simple-reflection
chains. Using the normalized completion instead gives the exact
finite identity $\mathsf b_mT_1T_2=\mathsf a_m$.

For $m=3$, the correction has the following form. On the appropriate
spherical domain, the $S_3$ identity
$C_{\{s,t\}}=C_sC_tC_s-C_s$ gives
\[
 \mathcal R_3=\mathcal{BS}_3(C_2C_3-1)(C_1C_2-1),
 \qquad C_i=C_{s_i}.
\]
For $m=2,3,4,5$ the uncorrected finite product
$\mathsf b_mL_m-\mathsf a_m$ has respectively $0,3,10,20$
distinct canonical terms. These are exact finite checks; the
canonical-tail and summand statements above hold for all $m$.
